\chapter{Introduction} 
\label{chap:introduction} 
It is a well-known fact, that Sweden is a country with big variations of the
daily hours of sun throughout the year. Far up north the sun does not set
anymore in January. Even in Stockholm in January, the darkest month of the
year, the sun rises at 8:47 am and sets at 2:55 pm
(cf. \cite{other:visitsweden}). According to \myemph{Sveriges Radio} "many Swedes
suffer from the winter blues or seasonal affective disorder" (cf.
\cite{other:sverigesradio}). In a strong winter every source of light is a
source of happiness. This is why a wakeup light, which is based on a strong
light source (at least 10 W RGB LED), is able to give one the optimal start into
a dark winter day with an artificial sunrise as bright as a real sun shining
through the window.  
\newpar
This report describes the Swakeup (from engl. ”Swedish Wakeup Light”), a device
communicating to the user not only through light. It does not simply wake one
up, but also gives one information about social media, latest mails, calendar
and weather. The user interface consits besides of a high-power LED of an OLED
screen. Swakeup is also part of the IoT as it has the ability to communicate
via IEEE 802.11. This of course enables a lot of possibilities e.g. connecting
your phone to the wakeup light. A lot of effort has been put into the designing
maxim, that everything should be as small as possible. The whole electronics
fit on an base area of 5 cm x 4 cm. So the Swakeup fits smoothly on the bedside
table. And honestly: What is the last thing people do before they go to sleep?
Right! They look on their phones. That is why Swakeup comes with a USB charger
for your phone as well, so you do not have to leave the bed anymore. Another
design maxim of this product is cheapness. Everybody should be able to buy one.
As all engineering work is open-source, it gives people (with the
corresponding knowledge) the opportunity to build a wakeup light themselves,
look up what this device is doing with their personal data or even contribute
to the product.

\begin{figure}[H]
		\centering
		\label{fig:block_simp}
		\includegraphics[width=0.2\textwidth]{./fig/github.jpg}
		%\caption{Simplified blockdiagram of Swakeup.}
\end{figure}
\begin{center}
\url{www.github.com/m3x1m0m/swakeup}
\end{center}
